% TEMPLATE-MILC-v3.tex - plantilla maestra UNICA de las guias MILC v3.
%
% La usan las cuatro familias de guias, cada una con su driver:
%   · Grados 6-11          -> scripts/build-guias-g11.py
%   · Bebras               -> scripts/build-guias-bebras.py
%   · Semillero            -> scripts/build-guias-semillero.py
%   · Territorio interior  -> scripts/build-guias-territorio-interior.py
%
% Antes habia tres copias casi identicas (~400 lineas iguales, 6 distintas):
% cada mejora habia que aplicarla tres veces, y en la practica no se hacia
% (el motor de recursos vivia solo en dos de ellas). Lo que varia entre
% familias esta parametrizado en 6 placeholders que cada driver llena
% (se nombran aqui SIN su sintaxis real: el builder reemplaza por texto plano,
% tambien dentro de los comentarios, y un valor multilinea rompe el preambulo):
%
%   HEADER_IZQ        encabezado izquierdo de pagina
%   HEADER_DER        encabezado derecho de pagina
%   PORTADA_ETIQUETA  rotulo sobre el numero grande (GRADO/MOMENTO/MODULO)
%   PORTADA_SERIE     linea de serie de la portada
%   PORTADA_ID        identificador de la guia en la portada
%   PORTADA_CIRCULO   el circulito de grado (°), o vacio si no aplica
%   PDF_TITULO        titulo en texto plano (sin \textbf ni comillas TeX) para
%                     los metadatos del PDF (/Title); lo produce lib_xelatex.texto_pdf
%
% Composicion (auditoria 2026-09-03): las bandas de estacion piden espacio
% (\Needspace*) y prohiben el salto justo despues (\nopagebreak), asi no quedan
% huerfanas al pie; las cajas largas (softbox, drawbox, infoband, puentebox) son
% `breakable`, asi una caja de media pagina ya no arrastra todo a la siguiente
% dejando la anterior al 40 %. Todo texto sobre mostaza o verde va en milcNegro
% (blanco sobre #E5B400 da 1,93:1 y sobre #58A923 2,95:1; ninguno pasa WCAG AA).
% El resto de claves las documenta content/guias/_SCHEMA.md. El script valida
% que no queden placeholders sin reemplazar antes de escribir el archivo final.

% Etiquetado PDF (latex-lab): /Lang, estructura y texto alternativo de las
% figuras, para lectores de pantalla. Debe ir ANTES de \documentclass.
\DocumentMetadata{lang=es-CO,tagging=on}
\documentclass[11pt,letterpaper]{article}
% headheight=24pt: la cabecera derecha lleva dos lineas (identidad de la guia
% y estacion actual). headsep se acorta en la misma medida para que el bloque
% de texto quede exactamente donde estaba con headheight=12pt.
\usepackage[margin=2.0cm,headheight=24pt,headsep=13pt]{geometry}
\usepackage[table]{xcolor}
\usepackage{fontspec}
\usepackage[spanish]{babel}
\usepackage{tikz}
% Librerías TikZ para los diagramas de las guías (bloque `recursos:`):
%  · babel  → babel en español vuelve activos caracteres como «>», y TikZ los usa
%             en sus opciones (>=stealth, ->). Sin esta librería, cualquier
%             diagrama con flechas revienta con «\language@active@arg> has an
%             extra }». Debe ir ANTES de que se dibuje nada.
%  · positioning        → sintaxis «below left=2pt and 3pt».
%  · decorations.pathreplacing → llaves ({brace}) para señalar rangos.
\usetikzlibrary{babel,positioning,decorations.pathreplacing}
\usepackage{graphicx}
% Circuitos: el trabajo de electrónica se hace hoy con micro:bit y sus sensores
% integrados (MakeCode, sin cableado), así que no se necesitan esquemáticos.
% Si algún día entran montajes con protoboard, basta descomentar esta línea:
% los diagramas .tex/.tikz declarados en `recursos:` ya se incrustan solos.
% \usepackage{circuitikz}
\usepackage[most]{tcolorbox}
\usepackage{tabularx}
\usepackage{array}
\usepackage{fancyhdr}
\usepackage{titlesec}
\usepackage[document]{ragged2e}  % bandera a la derecha en todo el cuerpo
\usepackage{enumitem}
\usepackage{microtype}
\usepackage{etoolbox}
\usepackage{needspace}
% hyperref al final del preambulo: metadatos del PDF (titulo, autor, idioma,
% familia), marcadores por estacion (\pdfbookmark) y enlaces sin recuadro.
\usepackage[hidelinks,
  pdftitle={Narrativa oral — tres minutos con voz propia},
  pdfauthor={PhD. Álvaro Cárdenas Orozco},
  pdfsubject={Tecnología e Informática · Grado 8 · Guía 23 · MILC v3},
  pdflang={es-CO},
  bookmarksopen=true,bookmarksnumbered=false]{hyperref}

% Helvetica Neue no trae la flecha U+2192, asi que cada «→» escrito literalmente
% en el YAML se compilaba a NADA, en silencio (462 flechas en 61 guias: rutas del
% tipo «paleta → tipografia → lockup» salian rotas). Mapeamos el caracter a su
% version matematica, que si tiene glifo. Asi el autor puede seguir escribiendo
% «→» comodamente en el YAML.
\usepackage{newunicodechar}
\newunicodechar{→}{\ensuremath{\rightarrow}}
% Mismo problema con los simbolos de marca y vinetas: el «✓» de «una palomita ✓»
% salia como cuadrito vacio en el PDF de todas las guias que lo usaban.
\usepackage{pifont}
\newunicodechar{✓}{\ding{51}}
\newunicodechar{✗}{\ding{55}}
\newunicodechar{✔}{\ding{52}}
\newunicodechar{•}{\textbullet}
\newunicodechar{≥}{\ensuremath{\geq}}
\newunicodechar{≤}{\ensuremath{\leq}}

\setmainfont{Helvetica Neue}
\setsansfont{Helvetica Neue}
\setlength{\parindent}{0pt}
\setlength{\parskip}{6pt}
% Sin particion de palabras: en 6.º-7.º una palabra cortada por guion al final
% de linea («Comuni-cacion») frena la lectura mucho mas que una linea corta.
\hyphenpenalty=10000
\exhyphenpenalty=10000
\pretolerance=2000
\tolerance=3000
\setlength{\emergencystretch}{3em}
\linespread{1.10}
\renewcommand{\arraystretch}{1.25}
\setlist{leftmargin=5mm,itemsep=1mm,topsep=1mm}
\newcolumntype{Y}{>{\RaggedRight\arraybackslash}X}

\definecolor{milcMagenta}{HTML}{E6007E}
\definecolor{milcVino}{HTML}{5A0038}
\definecolor{milcTurquesa}{HTML}{0093A5}
\definecolor{milcVerde}{HTML}{58A923}
\definecolor{milcMostaza}{HTML}{E5B400}
\definecolor{milcOcre}{HTML}{B86600}
\definecolor{milcNegro}{HTML}{241816}
\definecolor{milcGris}{HTML}{F7F7F4}
\definecolor{milcLinea}{HTML}{E8E2DD}
\definecolor{milcGradePrimary}{HTML}{FF6600}
\definecolor{milcGradeDark}{HTML}{9A3E00}
\definecolor{milcGradeSoft}{HTML}{FFE4D0}
\definecolor{milcGradeAccent}{HTML}{CFE7FF}
\definecolor{milcGradeText}{HTML}{FFFFFF}
\color{milcNegro}

\pagestyle{fancy}
\fancyhf{}
% Cabecera de dos lineas: identidad de la guia arriba y, a la derecha y en
% gris, la estacion en curso (\markright desde \milcestacion). Antes de la
% primera estacion la segunda linea va vacia; el \strut mantiene la altura
% para que la linea de arriba no baile entre paginas.
\lhead{\small Tecnología e Informática · Grado 8\\[-1pt]{\footnotesize\strut}}
\rhead{\small Guía 23 · MILC v3\\[-1pt]{\footnotesize\strut\color{milcNegro!65}\rightmark}}
\cfoot{\small\thepage}
\renewcommand{\headrulewidth}{0pt}

% Sobre mostaza (#E5B400) y verde (#58A923) el texto va en milcNegro; sobre el
% resto de la paleta, en blanco. Se usa dentro de fonttitle/fontupper, donde un
% \color posterior manda sobre coltitle.
\newcommand{\milctintasobre}[1]{%
  \ifboolexpr{test{\ifstrequal{#1}{milcMostaza}} or test{\ifstrequal{#1}{milcVerde}}}%
    {\color{milcNegro}}{\color{white}}}

\newtcolorbox{titlebox}[1]{
  enhanced,colback=#1,colframe=#1,arc=7mm,boxrule=0pt,
  left=7mm,right=7mm,top=3mm,bottom=3mm,
  before skip=8pt,after skip=8pt,
  fontupper=\sffamily\bfseries\Large\color{white}
}

\newtcolorbox{softbox}[3]{
  enhanced,breakable,pad at break=3mm,
  colback=#2,colframe=#1,borderline west={4pt}{0pt}{#1},
  arc=2mm,boxrule=.4pt,left=4mm,right=4mm,top=3mm,bottom=3mm,
  before skip=6pt,after skip=8pt,title={#3},coltitle=white,
  colbacktitle=#1,fonttitle=\sffamily\bfseries\milctintasobre{#1},fontupper=\RaggedRight,
  attach boxed title to top left={xshift=3mm,yshift=-2mm},
  boxed title style={arc=2mm, boxrule=0pt}
}

\newtcolorbox{drawbox}[2]{
  enhanced,breakable,pad at break=4mm,
  colback=white,colframe=#1,arc=4mm,boxrule=1pt,
  left=5mm,right=5mm,top=4mm,bottom=4mm,before skip=8pt,after skip=8pt,
  title={#2},coltitle=white,colbacktitle=#1,fonttitle=\sffamily\bfseries\milctintasobre{#1},
  fontupper=\RaggedRight,
  attach boxed title to top left={xshift=4mm,yshift=-2mm},
  boxed title style={arc=3mm, boxrule=0pt}
}

\newtcolorbox{infoband}[1]{
  enhanced,breakable,pad at break=2mm,colback=#1!8,colframe=#1!50,arc=2mm,boxrule=.5pt,
  left=4mm,right=4mm,top=2mm,bottom=2mm,before skip=4pt,after skip=4pt,
  fontupper=\small\RaggedRight
}

% Puente narrativo entre secciones (contrato editorial Conéctate).
% Da continuidad: "Acabas de X. Ahora vas a Y porque Z."
% Visualmente: banda izquierda en color del grado, texto en itálica pequeña.
\newtcolorbox{puentebox}{
  enhanced,breakable,pad at break=2mm,colback=milcGradeSoft,colframe=milcGradeDark,
  borderline west={3pt}{0pt}{milcGradeDark},
  arc=0mm,boxrule=0pt,left=5mm,right=4mm,top=2.5mm,bottom=2.5mm,
  before skip=4pt,after skip=8pt,
  fontupper=\itshape\small\color{milcNegro}\RaggedRight
}
% La flecha va en modo matematico a proposito: Helvetica Neue NO tiene el glifo
% U+2192, asi que \textrightarrow se compilaba a NADA («Missing character: There
% is no → in font Helvetica Neue») y el puente salia sin flecha. En modo
% matematico la flecha viene de la fuente matematica y si se dibuja.
\newcommand{\puente}[1]{%
  \begin{puentebox}%
    \textcolor{milcGradeDark}{$\rightarrow$}\hspace{2mm}#1%
  \end{puentebox}%
}

\newcommand{\linead}[1]{\noindent\color{milcLinea}\rule{#1}{0.5pt}\color{milcNegro}}
\newcommand{\respuesta}[1]{\par\vspace{2mm}\foreach \n in {1,...,#1}{\noindent\color{milcLinea}\rule{\linewidth}{0.5pt}\par\vspace{5mm}}\color{milcNegro}}
\newcommand{\checkbox}{\textcolor{milcGradeDark}{\fbox{\phantom{x}}}\hspace{5pt}}

% ─── Listas aireadas para las guías socioemocionales ────────────────────
% Componer las verificaciones como «\checkbox texto\\» deja el texto que salta
% de línea debajo del cuadrito y apelmaza el bloque. Estos entornos alinean la
% continuación y airean los ítems. Los usa Territorio Interior; cualquier
% familia puede hacerlo.
\newlist{checklista}{itemize}{1}
\setlist[checklista]{
  label=\textcolor{milcGradeDark}{\fbox{\phantom{x}}},
  leftmargin=9mm, labelsep=3.2mm, itemsep=2.4mm,
  topsep=2.5mm, parsep=0pt, partopsep=0pt, before=\RaggedRight
}
\newlist{pasoslista}{enumerate}{1}
\setlist[pasoslista]{
  label=\textcolor{milcGradeDark}{\sffamily\bfseries\arabic*.},
  leftmargin=8mm, labelsep=3mm, itemsep=2.8mm,
  topsep=1mm, parsep=0pt, partopsep=0pt, before=\RaggedRight
}

\newcommand{\phasecard}[4]{
\begin{tcolorbox}[enhanced,colback=#3,colframe=#2,arc=3mm,boxrule=.5pt,height=3.05cm,valign=center,halign=center,left=1.5mm,right=1.5mm,top=2mm,bottom=2mm]
{\sffamily\bfseries\fontsize{22}{24}\selectfont\textcolor{#2}{#1}}\par
{\sffamily\bfseries\small #4}
\end{tcolorbox}
}

% ── Piezas del rediseno 2026 (legibilidad 6.º-11.º) ───────────────────
% Banda de estacion: reemplaza el titlebox macizo. Titulo a la izquierda,
% dato practico (tiempo/modalidad) a la derecha, en una sola linea.
% Nunca huerfana: pide 9 lineas libres (o salta de pagina rellenando con
% \vfill) y prohibe el salto justo despues de la banda. Nueve y no seis:
% la caja partible que sigue necesita ~80pt (titulo adosado, relleno y dos
% lineas) para arrancar en la misma pagina; con menos, tcolorbox la manda
% entera a la siguiente y la banda se queda sola. El penalty 10000 va
% en `after`, ANTES del espacio posterior: un `after skip` (aunque sea 0pt)
% es un \vskip tras la caja, y ese glue es un punto de corte legal que un
% \nopagebreak escrito despues de \end{tcolorbox} ya no cierra. Cada banda
% deja marcador PDF y actualiza la cabecera derecha.
\newcounter{milcestacion}
\newcommand{\milcestacion}[2]{%
\Needspace*{9\baselineskip}%
\stepcounter{milcestacion}%
\pdfbookmark[1]{#2}{milcestacion\themilcestacion}%
\markright{#2}%
\begin{tcolorbox}[enhanced,colback=#1,colframe=#1,arc=2mm,boxrule=0pt,
  left=5mm,right=5mm,top=2.6mm,bottom=2.6mm,before skip=11pt,
  after={\par\nopagebreak\vskip7pt}]
{\sffamily\bfseries\fontsize{15}{18}\selectfont\milctintasobre{#1}#2}
\end{tcolorbox}}

% Tarjeta de paso numerada: sustituye las tablas «Pilar / Aplicacion», que
% eran formato de consulta docente, no de estudiante. El numero va en un
% circulo sobre el borde para que la secuencia se lea de un vistazo.
\newcommand{\milcpaso}[3]{%
\Needspace{4\baselineskip}%
\begin{tcolorbox}[enhanced,colback=white,colframe=milcLinea,arc=1.5mm,boxrule=.9pt,
  left=14mm,right=4mm,top=3mm,bottom=3mm,before skip=4pt,after skip=4pt,
  fontupper=\RaggedRight,
  overlay={\node[anchor=north west,circle,fill=#2,text=white,inner sep=0pt,
    minimum size=7.4mm,font=\sffamily\bfseries\small]
    at ([xshift=3.6mm,yshift=-3.2mm]frame.north west) {#1};}]
#3
\end{tcolorbox}}

% Casilla marcable de tamano util con lapiz (la anterior era \fbox{\phantom{x}},
% de alto variable segun el texto que la seguia).
\newcommand{\milcmarca}{\framebox[3.8mm]{\rule{0pt}{3.4mm}}}

% Fila marcable de lista suelta (checks de la estacion 1).
\newcommand{\milccheckrow}[1]{%
\par\vspace{1.8mm}\noindent
\makebox[6.4mm][l]{\raisebox{-.55mm}{\textcolor{milcNegro}{\milcmarca}}}%
\parbox[t]{\dimexpr\linewidth-6.4mm\relax}{\RaggedRight #1}\par}

% Celda marcable dentro de la rubrica (tabla de autoevaluacion). Misma
% sangria francesa, pero pensada para vivir dentro de una columna X.
\newcommand{\milcopcion}[1]{%
\makebox[5.2mm][l]{\raisebox{-.55mm}{\textcolor{milcNegro}{\milcmarca}}}%
\parbox[t]{\dimexpr\linewidth-5.2mm\relax}{\RaggedRight #1}}

% ── Recursos visuales de la guía ──────────────────────────────────────
% Declarados en el bloque `recursos:` del YAML y emitidos por el builder.
% \guiaFigura   → imagen rasterizada o PDF (foto, carta celeste, montaje).
% \guiaDiagrama → fuente TikZ/circuitikz (.tex) incrustada como vector:
%                 es la vía para circuitos y esquemáticos editables.
% [#1] = texto alternativo (alt) · #2 = ruta relativa al .tex · #3 = pie
% El alt llega al PDF etiquetado (/Alt) y lo lee el lector de pantalla.
\newcommand{\guiaFigura}[3][]{%
\begin{center}
\includegraphics[width=0.88\linewidth,height=0.38\textheight,keepaspectratio,alt={#1}]{#2}\\[1.5mm]
{\footnotesize\itshape\textcolor{milcNegro!75}{#3}}
\end{center}
}

\newcommand{\guiaDiagrama}[2]{%
\begin{center}
\input{#1}\\[1.5mm]
{\footnotesize\itshape\textcolor{milcNegro!75}{#2}}
\end{center}
}

\begin{document}
\thispagestyle{empty}

% ── Portada ───────────────────────────────────────────────────────────
% Antes: un rectangulo de color a pagina completa. Se veia bien en pantalla,
% pero en la fotocopiadora del colegio se comia el toner y salia gris sucio.
% Ahora la banda ocupa el tercio superior y el resto va en blanco.
\begin{tikzpicture}[remember picture,overlay]
  \path[fill=milcGradePrimary]
    (current page.north west) rectangle ([yshift=-10.6cm]current page.north east);

  \node[anchor=north west,text=milcGradeText] at ([xshift=2.0cm,yshift=-1.55cm]current page.north west)
    {\sffamily\bfseries\fontsize{12}{14}\selectfont GRADO};
  \node[anchor=north west,text=milcGradeText,opacity=.85] at ([xshift=2.0cm,yshift=-2.35cm]current page.north west)
    {\sffamily\bfseries\fontsize{8.5}{10}\selectfont SERIE GUÍAS · TECNOLOGÍA E INFORMÁTICA};
  \node[anchor=north east,text=milcGradeText,opacity=.85,align=right] at ([xshift=-2.0cm,yshift=-1.55cm]current page.north east)
    {\sffamily\bfseries\fontsize{9}{12}\selectfont I.E. Sor María Juliana\\Cartago, Valle del Cauca};

  \node[anchor=west,text=milcGradeText] (gradonum) at ([xshift=1.85cm,yshift=-7.1cm]current page.north west)
    {\sffamily\bfseries\fontsize{112}{112}\selectfont 8};
  % El circulito de grado (°) solo aplica a las guias de grado («11°»); en
  % Bebras y Semillero el numero grande es un momento/modulo, no un grado.
  % Se posiciona relativo al nodo `gradonum`, no en coordenadas absolutas.
  \draw[milcGradeText,line width=1.6pt]
    ([xshift=2.0mm,yshift=-4.5mm]gradonum.north east) circle (.15cm);

  \node[anchor=west,text=milcGradeText,text width=11.4cm,align=left] at ([xshift=1.5cm,yshift=0.5cm]gradonum.east)
    {
     {\sffamily\bfseries\fontsize{9}{11}\selectfont\MakeUppercase{Período 3 · Multimedia, narrativa y ciberseguridad}}\\[3mm]
     {\sffamily\bfseries\fontsize{21}{25}\selectfont\setlength{\baselineskip}{27pt}Tres minutos\\[4pt]con voz propia}\\[3.5mm]
     {\sffamily\bfseries\fontsize{15}{18}\selectfont Guía 23-8-TIC}
    };
\end{tikzpicture}

\vspace*{9.4cm}
\vfill

{\sffamily\bfseries\fontsize{8}{10}\selectfont\textcolor{milcGradeDark}{LO QUE TE LLEVAS HOY}}

\begin{tcolorbox}[enhanced,colback=white,colframe=milcNegro,arc=2mm,boxrule=1.6pt,
  left=5mm,right=5mm,top=4mm,bottom=4mm,before skip=3pt,after skip=12pt,
  fontupper=\RaggedRight\fontsize{11}{15.5}\selectfont]
Un discurso oral de tres minutos sobre algo que te importa, que aprendiste o que quisieras cambiar, con guion escrito de 400 a 450 palabras en tres actos: gancho de 30 segundos, desarrollo de dos minutos y cierre de 30 segundos con una petición concreta. Grabado con el celular en una sola toma o con máximo tres cortes, y una autocrítica con dos cosas que funcionaron y una que cambiarías.
\end{tcolorbox}

\begin{tcolorbox}[enhanced,colback=milcGris,colframe=milcGris,arc=2mm,boxrule=0pt,
  left=5mm,right=5mm,top=3.5mm,bottom=3.5mm,after skip=12pt]
{\sffamily\bfseries\fontsize{8}{10}\selectfont\textcolor{milcNegro!55}{ESTA GUÍA ES DE}}\\[3mm]
\begin{tabularx}{\linewidth}{X p{3.2cm} p{3.2cm}}
\linead{\linewidth} & \linead{3.0cm} & \linead{3.0cm}\\[-1mm]
{\fontsize{8.5}{10}\selectfont\textcolor{milcNegro!55}{Nombre completo}} &
{\fontsize{8.5}{10}\selectfont\textcolor{milcNegro!55}{Grupo}} &
{\fontsize{8.5}{10}\selectfont\textcolor{milcNegro!55}{Fecha}}\\
\end{tabularx}
\end{tcolorbox}

\vfill

\noindent\textcolor{milcLinea}{\rule{\linewidth}{1pt}}\\[2mm]
{\sffamily\bfseries\fontsize{9.5}{12}\selectfont Autor: Dr. Álvaro Cárdenas Orozco}\\[1mm]
{\sffamily\fontsize{8.5}{11}\selectfont\textcolor{milcNegro!55}{Metodología MILC · apertura ancestral · pensamiento computacional · triángulo Dussel–estoicismo–Floridi}}

\newpage

\section*{Apertura: Narrativa oral --- tres minutos con voz propia}

\begin{softbox}{milcGradeDark}{milcGradeSoft}{Saber ancestral}
Desde el cerro Munchique, en Santander de Quilichao, una emisora del pueblo nasa abre cada mañana con la misma frase: «Muy buenos días, pa'yumat. Gracias por dejarnos entrar» (Aley, 2012). \textbf{Pa'yumat} quiere decir bienvenido, y también pedir permiso para entrar. Así se llama la emisora, y hace parte del Tejido de Comunicación de la Asociación de Cabildos Indígenas del Norte del Cauca. Su razón de existir cabe en una línea del propio Tejido: «recuperar la palabra y proyectarla, para que otros no hablen por nosotros». Los programas se hacen en nasa yuwe y en castellano. Los hacen comunicadores de la misma comunidad, que cuentan lo que pasa en los resguardos con su propia voz. Fíjate en el saludo: no empieza con «hoy vamos a hablar de». Pide permiso, agradece, y entra. La \textbf{cara de exclusión}: en 2012 hubo atentados contra las antenas y torres de la zona; la señal perdió alcance y los equipos tuvieron que bajarse del cerro. Hablar con voz propia, allá, tiene un costo que aquí no tiene. Hoy vas a hablar tres minutos de algo que te importa, con tu voz, y a empezar como Pa'yumat: pidiendo permiso, no anunciando.\par\smallskip{\footnotesize\textcolor{milcNegro!70}{\textbf{Fuente.} Tejido de Comunicación ACIN. (s.\,f.). \emph{Radio Pa'yumat}. Consultado el 7 de septiembre de 2026. · Aley, P. (2012, 22 de septiembre). Conozca la emisora de la resistencia indígena. \emph{El Tiempo}.}}
\end{softbox}

\begin{softbox}{milcTurquesa}{milcGris}{Saber contemporáneo}
Un discurso corto se sostiene en \textbf{tres actos}. El \textbf{gancho}, unos 30 segundos: la primera frase no anuncia, entra. Puede ser una escena, una pregunta que incomoda, un dato con fuente o una imagen concreta. El \textbf{desarrollo}, unos dos minutos: aquí va la historia con algo en juego, una tensión que no se resuelve de inmediato. El \textbf{cierre}, unos 30 segundos: vuelve al gancho y pide algo concreto a quien escucha. Tres minutos hablados son entre 400 y 450 palabras a un ritmo normal; por eso el guion se escribe y se cuenta con cronómetro. Se escribe en prosa, como se habla, no en viñetas. Y no se lee: leer mata el discurso aunque el texto sea bueno. Se ensaya hasta que salga contado. Los errores más comunes son cinco: empezar con un saludo largo, leer, contar sin tensión, cerrar con un «gracias» y pasarse del tiempo.
\end{softbox}

\begin{softbox}{milcMostaza}{milcGris}{Pregunta-puente}
\textit{Radio Pa'yumat entra pidiendo permiso, no anunciando el programa. Si tu discurso empieza con «hoy les voy a hablar de», ¿cuántos de los 30 segundos del gancho gastaste sin decir nada?}
\end{softbox}

\begin{softbox}{milcVerde}{milcGris}{Saber y saber hacer}
Al terminar podrás: (1) \textbf{identificar} los tres actos en un cuento o una película que conoces; (2) \textbf{crear} un guion de 400 a 450 palabras con gancho, desarrollo y cierre; (3) \textbf{evaluar} tu grabación con dos aciertos y un cambio; (4) \textbf{explicar} por qué un gancho pide permiso en vez de anunciar.
\end{softbox}



\small
\begin{tabularx}{\textwidth}{>{\bfseries}p{3.0cm}Y}
\rowcolor{milcGradePrimary!12}Autor & Dr. Álvaro Cárdenas Orozco\\
DBA & Producción multimedia ética y comportamiento responsable en entornos digitales (MEN, T\&I)\\
Referentes & Dussel (estética de la liberación) · Ley 1336 · Floridi (infoesfera)\\
Producto final & Un discurso oral de tres minutos sobre algo que te importa, que aprendiste o que quisieras cambiar, con guion escrito de 400 a 450 palabras en tres actos: gancho de 30 segundos, desarrollo de dos minutos y cierre de 30 segundos con una petición concreta. Grabado con el celular en una sola toma o con máximo tres cortes, y una autocrítica con dos cosas que funcionaron y una que cambiarías.\\
\end{tabularx}
\normalsize

\puente{En la sesión 2 armaste caminos con un menú que lleva la esencia. Hoy cuentas algo con tu voz en tres minutos, y el gancho hace lo que hacía el menú. En la sesión 4 vas a editar imágenes con ética.}

\milcestacion{milcGradeDark}{Ruta MILC: así vamos a aprender}
\begin{tabularx}{\textwidth}{>{\centering\arraybackslash}X>{\centering\arraybackslash}X>{\centering\arraybackslash}X>{\centering\arraybackslash}X}
\phasecard{1}{milcVerde}{milcGris}{Escucha\\[-1mm]\small Observo, escucho y pregunto.} &
\phasecard{2}{milcTurquesa}{milcGris}{Entender\\[-1mm]\small Sistematizo y ordeno.} &
\phasecard{3}{milcMagenta}{milcGris}{Hacer\\[-1mm]\small Construyo y aplico.} &
\phasecard{4}{milcVino}{milcGris}{Cierre\\[-1mm]\small Evalúo y reflexiono.}
\end{tabularx}


\puente{Antes de escribir, vas a encontrar los tres actos en una historia que ya conoces. Si sabes dónde te atrapó y dónde te soltó, sabes qué es un gancho y un cierre.}

\milcestacion{milcVerde}{Estación 1 de 3 · Escucha}

\begin{softbox}{milcVerde}{milcGris}{Escena para escuchar}
\textbf{Actividad 1 · IDENTIFICA --- Los tres actos de una historia que conoces} (15 min · individual). (1) Elige una historia que recuerdes bien: una película, un capítulo de serie, un cuento de tu abuela, un partido. (2) Escribe cómo empezó y qué te atrapó en el primer minuto: ese es el gancho. (3) Escribe qué estaba en juego en la mitad y en qué momento no sabías cómo iba a terminar: ese es el desarrollo. (4) Escribe cómo cerró y si el final volvió a algo del principio: ese es el cierre. (5) Anota en una línea qué habrías cambiado del gancho o del cierre.
\end{softbox}

{\sffamily\bfseries\fontsize{8}{10}\selectfont\textcolor{milcNegro!55}{MARCA CUANDO LO HAYAS HECHO}}
\milccheckrow{Tengo el gancho, el desarrollo y el cierre de una historia que conozco.}
\milccheckrow{Escribí qué estaba en juego en la mitad.}
\milccheckrow{Anoté si el cierre volvió al principio.}

\begin{infoband}{milcVerde}


\textbf{Lo que va al cuaderno (Actividad 1 · IDENTIFICA).} \textbf{Título:} «Actividad 1 --- Los tres actos de una historia que conoces». \textbf{Formato:} tres bloques de dos o tres líneas, gancho, desarrollo y cierre, y la línea de lo que cambiarías. \textbf{Extensión:} media página. \textbf{Sabes que terminaste cuando} puedes decir en qué momento exacto la historia te atrapó y en cuál te soltó.
\end{infoband}

\puente{Ya viste los tres actos en otra historia. Ahora aprendes los cinco ganchos y, con tu pareja, escribes y cronometras tu guion.}

\milcestacion{milcTurquesa}{Estación 2 de 3 · Entender}

\begin{softbox}{milcTurquesa}{milcGris}{Lo que vas a entender}
\textbf{Actividad 2 · CREA --- El guion de tres minutos} (30 min · parejas). (1) Elige el tema: algo que te importa, algo que aprendiste, algo que quisieras cambiar del colegio o del barrio. (2) Escribe el gancho en tres o cuatro frases, sin saludo largo y sin «hoy voy a hablar de»: una escena, una pregunta, un dato con fuente o una imagen. (3) Escribe el desarrollo en prosa, con algo en juego, hasta llegar a unas 300 palabras. (4) Escribe el cierre: vuelve al gancho y pide algo concreto a quien escucha. (5) Léelo a tu pareja con cronómetro; si pasa de 3:15, recorta; si no llega a 2:45, falta desarrollo.
\end{softbox}

\milcpaso{1}{milcTurquesa}{\textbf{Los cinco ganchos.} Una escena personal corta: «cuando tenía diez años…». Una pregunta que incomoda: «¿cuántas horas de tu vida real llevas en el celular?». Un dato con fuente: solo si sabes de dónde salió. Una imagen concreta: algo que se pueda ver. Un permiso, como Pa'yumat: «gracias por escucharme tres minutos».}
\milcpaso{2}{milcTurquesa}{\textbf{El desarrollo con algo en juego.} Una lista de datos no es un desarrollo. Hace falta una tensión: algo que podía salir mal, una decisión difícil, una pregunta sin respuesta hasta el final. Si quien escucha ya sabe cómo termina, deja de escuchar.}
\milcpaso{3}{milcTurquesa}{\textbf{El cierre que vuelve.} El mejor cierre toca algo del gancho: la misma escena, la misma pregunta ya respondida. Y pide algo concreto: «pregúntale a tu abuela», «apaga el celular una hora». Un «gracias» no es un cierre.}
\milcpaso{4}{milcTurquesa}{\textbf{Cómo se hace, paso a paso.} (1) Elegir el tema. (2) Escribir el gancho. (3) Escribir el desarrollo con tensión. (4) Escribir el cierre que vuelve y pide. (5) Cronometrar con tu pareja. (6) Recortar o ampliar.}

\begin{drawbox}{milcTurquesa}{El discurso de tres minutos, en una mirada}
\textbf{Gancho (30 s):} entra, no anuncia. \\[1mm] \textbf{Desarrollo (2 min):} algo en juego. \\[1mm] \textbf{Cierre (30 s):} vuelve al gancho y pide algo. \\[1mm] \textbf{Guion:} 400 a 450 palabras en prosa. \\[1mm] \textbf{Regla:} se cuenta, no se lee.
\end{drawbox}

\begin{softbox}{milcMostaza}{milcGris}{Errores comunes y su corrección}
(1) \textbf{Saludo largo}: «buenos días, mi nombre es, hoy voy a hablar de» gasta el gancho. (2) \textbf{Leer}: aunque el texto sea bueno, suena a lectura. (3) \textbf{Sin tensión}: una lista de datos que no va a ningún lado. (4) \textbf{Cierre con «gracias»}: no pide nada ni vuelve al principio. (5) \textbf{Pasarse del tiempo}: a los cuatro minutos ya nadie escucha.
\end{softbox}

\begin{infoband}{milcTurquesa}


\textbf{Lo que va al cuaderno (Actividad 2 · CREA).} \textbf{Título:} «Actividad 2 --- El guion de tres minutos». \textbf{Formato:} el guion completo en prosa con los tres actos marcados al margen, el conteo de palabras y el tiempo que marcó el cronómetro. \textbf{Extensión:} una página. \textbf{Sabes que terminaste cuando} el guion tiene entre 400 y 450 palabras, la primera frase no anuncia y el cierre pide algo.
\end{infoband}

\puente{Ya tienes el guion. Toca ensayar sin leer, grabar con el celular y hacer la autocrítica.}

\milcestacion{milcMagenta}{Estación 3 de 3 · Hacer}

\begin{softbox}{milcMagenta}{milcGris}{Cómo vas a construir tu producto}
\textbf{Actividad 3 · EVALÚA --- Grabar y criticarse} (25 min · individual). (1) Ensaya dos veces sin mirar el guion; puedes tener tres palabras clave en un papel. (2) Graba con el celular en una sola toma; si te equivocas, sigue, y solo si es grave corta, máximo tres cortes. (3) Escucha la grabación completa con audífonos. (4) Escribe la autocrítica: dos cosas que funcionaron y una que cambiarías, con el minuto y segundo donde pasa. (5) Pásale a tu pareja solo los primeros 30 segundos y pregúntale si quiere seguir oyendo.
\end{softbox}

\milcpaso{1}{milcMagenta}{\textbf{Tu entrega tiene tres piezas.} (1) El guion con los tres actos marcados. (2) La grabación compartida con enlace o archivo. (3) La autocrítica con dos aciertos y un cambio, cada uno con su minuto y segundo.}
\milcpaso{2}{milcMagenta}{\textbf{Contar, no leer.} Tres palabras clave en un papel son suficientes. Si se te olvida una frase, di la idea con otras palabras. Nadie tiene el guion; nadie sabe que cambiaste algo.}
\milcpaso{3}{milcMagenta}{\textbf{La prueba de los 30 segundos.} Tu pareja oye el gancho y nada más. Si dice «quiero saber qué pasa», el gancho entró. Si dice «¿y de qué es?», anunciaste en vez de entrar.}
\milcpaso{4}{milcMagenta}{\textbf{La autocrítica con minuto y segundo.} «Estuve nervioso» no sirve. «En 1:40 aceleré y no se entendió la pregunta» sí sirve, porque dice qué ensayar la próxima vez.}

\begin{drawbox}{milcMagenta}{Antes de entregar}
\checkbox Guion de 400 a 450 palabras en prosa con tres actos marcados. \\[1mm] \checkbox La primera frase entra sin saludo largo ni «hoy voy a hablar de». \\[1mm] \checkbox El desarrollo tiene algo en juego. \\[1mm] \checkbox El cierre vuelve al gancho y pide algo concreto. \\[1mm] \checkbox Grabación de 2:45 a 3:15, una toma o máximo tres cortes, sin leer. \\[1mm] \checkbox Autocrítica con dos aciertos y un cambio, con minuto y segundo. \\[1mm] \checkbox Mi pareja oyó los primeros 30 segundos y dijo si quería seguir.
\end{drawbox}

\begin{softbox}{milcMagenta}{milcGris}{Plantilla de trabajo}
\textbf{Gancho.} \textit{La primera frase es …} \\[1mm] \textbf{En juego.} \textit{Lo que podía salir mal era …} \\[1mm] \textbf{Cierre.} \textit{Vuelvo a … y pido …} \\[1mm] \textbf{Tiempo.} \textit{… palabras, …:… en el cronómetro.} \\[1mm] \textbf{Cambio.} \textit{En …:… voy a …}
\end{softbox}

\begin{infoband}{milcMagenta}


\textbf{Lo que va al cuaderno (Actividad 3 · EVALÚA).} \textbf{Título:} «Actividad 3 --- Grabar y criticarse». \textbf{Formato:} el enlace o nombre del archivo de la grabación, la duración exacta, la autocrítica con dos aciertos y un cambio con su minuto y segundo, y lo que dijo tu pareja tras 30 segundos. \textbf{Extensión:} media página. \textbf{Sabes que terminaste cuando} la grabación dura entre 2:45 y 3:15 y tu cambio tiene un minuto y segundo.
\end{infoband}

\puente{Lo que entregas hoy: el guion, la grabación y la autocrítica. La prueba de calidad: tu pareja escucha los primeros 30 segundos y dice si quiere seguir oyendo o no.}

\milcestacion{milcMostaza}{Producto de la sesión}
\textbf{Discurso de tres minutos grabado + guion en tres actos + autocrítica con minuto y segundo + prueba de los 30 segundos}.
\begin{softbox}{milcMostaza}{milcGris}{Criterios de calidad}
(1) El guion tiene entre 400 y 450 palabras en prosa, con gancho, desarrollo y cierre marcados. (2) La primera frase entra sin saludo largo ni anuncio. (3) El desarrollo tiene algo en juego y el cierre vuelve al gancho y pide algo concreto. (4) La grabación dura entre 2:45 y 3:15, en una toma o máximo tres cortes, y no suena leída. (5) La autocrítica tiene dos aciertos y un cambio con minuto y segundo. (6) Tu pareja oyó los primeros 30 segundos y dijo si quería seguir.
\end{softbox}

\puente{Antes de cerrar, mira lo que hiciste desde cinco lados. Tres minutos contados con tu voz valen más que diez leídos de una pantalla.}

\milcestacion{milcVino}{Evaluación liberadora · cinco dimensiones}

\begin{softbox}{milcVino}{milcGris}{Sentido de la evaluación}
Evaluar no es solo revisar una nota. Es reconocer qué aprendí, qué cambió en mí, qué emociones manejé, cómo me relaciono con la ciudad y la comunidad, y qué le dejaré a quien viene detrás.
\end{softbox}

\textbf{Mi reflexión liberadora en cinco dimensiones.} Completa cada frase iniciada:

\small
\begin{tabularx}{\textwidth}{>{\bfseries\RaggedRight\arraybackslash}p{3.4cm}Y>{\RaggedRight\arraybackslash}p{4.6cm}}
\rowcolor{milcVino}\color{white}Dimensión & \color{white}\bfseries Mi respuesta (frame starter) & \color{white}\bfseries Algo que descubrí\\
1. Personal & Lo que más cambió en mí fue \linead{6.5cm} & \linead{4.2cm}\\
2. Emocional & La emoción más fuerte fue \linead{2.5cm} y la regulé \linead{3cm} & \linead{4.2cm}\\
3. Ciudadana & En democracia esto importa porque \linead{6cm} & \linead{4.2cm}\\
4. Local (Cartago) & En mi barrio sería útil para \linead{6.5cm} & \linead{4.2cm}\\
5. Intergeneracional & A mi abuela/abuelo le diría \linead{6.5cm} & \linead{4.2cm}\\
\end{tabularx}
\normalsize

\begin{infoband}{milcVino}
\textbf{Para la próxima sustentación.} Vuelve a esta tabla en seis meses. Verás que las respuestas cambian — no porque lo aprendido cambie, sino porque tú cambias. La evaluación liberadora es una conversación contigo a través del tiempo.
\end{infoband}


\puente{Para terminar, tres ideas sobre hablar con voz propia, contadas con palabras de esta guía: Dussel sobre la voz que irrumpe desde fuera, Marco Aurelio sobre meterse en el alma de quien habla, y la Onlife Initiative sobre lo que rechazamos porque no lo entendemos.}

\milcestacion{milcGradeDark}{Tres ideas para cerrar · Dussel, un estoico y Floridi}

\begin{softbox}{milcVino}{milcGris}{Enrique Dussel · lente del nosotros}
\textit{La conciencia ética es saber escuchar la voz del otro cuando irrumpe desde fuera de lo establecido.}

{\footnotesize\itshape\color{milcNegro!70}--- idea tomada de Enrique Dussel, contada con palabras de esta guía. No es una cita textual.}

Radio Pa'yumat existe para que otros no hablen por los nasa. Tu discurso es tu voz irrumpiendo con algo que te importa; y cuando escuches el de un compañero, la conciencia ética es oírlo aunque no sea lo que esperabas.

\textbf{¿A quién le hablan por encima en mi colegio, y qué diría con tres minutos propios?}
\end{softbox}
\linead{\linewidth}\\[1mm]
\linead{\linewidth}

\begin{softbox}{milcMostaza}{milcGris}{Estoicismo · Marco Aurelio · cuidado interior}
\textit{Acostúmbrate a atender de verdad a lo que dice el otro y a meterte, en lo posible, en su forma de ver.}

{\footnotesize\itshape\color{milcNegro!70}--- idea tomada de Marco Aurelio, contada con palabras de esta guía. No es una cita textual.}

Un gancho funciona cuando quien escucha se mete en tu escena. Y cuando tu pareja te dice «en 1:40 no entendí», está atendiendo de verdad. Recibir eso sin defenderse es lo que mejora el discurso.

\textbf{¿Qué discurso de un compañero escuché hoy de verdad, y qué me hizo ver distinto?}
\end{softbox}
\linead{\linewidth}\\[1mm]
\linead{\linewidth}

\begin{softbox}{milcTurquesa}{milcGris}{Luciano Floridi y la Onlife Initiative · lente de la infoesfera}
\textit{Tememos y rechazamos aquello a lo que no logramos darle sentido.}

{\footnotesize\itshape\color{milcNegro!70}--- idea tomada de Luciano Floridi y la Onlife Initiative, contada con palabras de esta guía. No es una cita textual.}

Un discurso que empieza anunciando no le da sentido a nada; el oyente se va. El gancho da sentido en 30 segundos: una escena, una pregunta. Lo que tiene sentido no se rechaza.

\textbf{¿Qué de mi tema me daba miedo contar, y qué pasó cuando le di una escena?}
\end{softbox}
\linead{\linewidth}\\[1mm]
\linead{\linewidth}

\begin{infoband}{milcNegro}
\textbf{Nota para el docente.} Las tres ideas están contadas con palabras de esta guía a partir del pensamiento de cada autor; no son citas textuales. De 9.º en adelante se trabajan con la obra y su referencia.
\end{infoband}

\milcestacion{milcVino}{Mi autoevaluación · marca una casilla por fila}

{\small
\setlength{\extrarowheight}{2.2pt}
\begin{tabularx}{\textwidth}{>{\RaggedRight\arraybackslash\bfseries}p{3.0cm}YYY}
\rowcolor{milcVino}\color{white}\bfseries Criterio & \color{white}\bfseries Lo logré & \color{white}\bfseries En proceso & \color{white}\bfseries Necesito apoyo\\[.6mm]
Comprensión del tema & \milcopcion{Explico las ideas con mis palabras.} & \milcopcion{Reconozco las ideas, falta precisión.} & \milcopcion{Necesito repasar lo básico.}\\[.8mm]
\rowcolor{milcGris}Pensamiento computacional & \milcopcion{Usé los pilares al armar mi producto.} & \milcopcion{Usé algunos pilares.} & \milcopcion{Necesito releer la estación 2.}\\[.8mm]
Aplicación y producto & \milcopcion{Mi producto cumple los criterios.} & \milcopcion{Está armado, con ajustes pendientes.} & \milcopcion{Aún está en construcción.}\\[.8mm]
\rowcolor{milcGris}Reflexión 5 dimensiones & \milcopcion{Respondí cada una con honestidad.} & \milcopcion{Respondí algunas con detalle.} & \milcopcion{Mis respuestas son muy generales.}\\[.8mm]
Triángulo de pensamiento & \milcopcion{Conecté las 3 voces con mi trabajo.} & \milcopcion{Conecté al menos una.} & \milcopcion{No las relaciono con mi proceso.}\\[.8mm]
\end{tabularx}}

\vspace{3mm}
\textbf{Hoy entendí que:}\\[1mm]
\linead{\linewidth}\\[1mm]
\linead{\linewidth}

\textbf{Mi compromiso para la próxima sesión es:}\\[1mm]
\linead{\linewidth}\\[1mm]
\linead{\linewidth}

\begin{softbox}{milcMostaza}{milcGris}{Cita de despedida · Marco Aurelio}
\textit{``El obstáculo en el camino se vuelve el camino. Lo que estorba al camino se vuelve camino.''}\\[1mm]
Cada error de hoy, bien observado, se convierte en pieza del camino que pisarás mañana.
\end{softbox}

\small
\begin{tabularx}{\textwidth}{>{\bfseries}p{3.6cm}Y}
\rowcolor{milcGradeDark!12}Nombre completo: & \linead{10cm}\\
Fecha: & \linead{4cm} \quad Curso/grupo: \linead{4cm}\\
Mi firma: & \linead{9cm}\\
\end{tabularx}
\normalsize

\begin{infoband}{milcGradeDark}
\textbf{Para la próxima sesión.} Llega con tu grabación y la autocrítica. En la sesión 4 vas a editar una imagen con recorte, brillo y contraste, y a responder cinco preguntas éticas antes de publicarla. La voz propia de hoy tiene una regla parecida: solo se usa la imagen de alguien con su permiso.
\end{infoband}

\end{document}
